Eine Flüssigkeitsmenge der Masse $m$ hat vor der Druckzunahme das Volumen $V_0$ und die Dichte $\rho_0 = m/V_0$. Bei der Druckzunahme $\Delta p$ ändert sich das Volumen um (FoTaBe S.~164)
\begin{align}
 \Delta V &= -\chi\,V_0\,\Delta p
\end{align}
Die Masse bleibt gleich, also ist die neue Dichte
\begin{align}
 \rho &= \frac{m}{V_0+\Delta V} = \frac{m}{V_0\left(1-\chi\,\Delta p\right)} = \frac{\rho_0}{1-\chi\,\Delta p}
\end{align}
Da $\chi\,\Delta p \ll 1$ ist (bei Flüssigkeiten und Drücken von einigen bar weniger als ein Tausendstel), gilt $\frac{1}{1-x}\approx 1+x$ und damit
\begin{align}
 \rho &\approx \rho_0\left(1+\chi\,\Delta p\right) \quad\Rightarrow\quad \result{\Delta\rho = \rho_0\,\chi\,\Delta p}
\end{align}
